% used for graphical images
\usepackage{graphicx} 

% control sizes of content
\usepackage{adjustbox}

% used for more color control
\usepackage{xcolor}

% used to control drawings on the axses
\usepackage{tikz}

% Setup for text margin
\usepackage[a4paper, margin=3cm]{geometry}

% Better chapters 
\usepackage{titlesec}

% For more fanzy heading
\usepackage{fancyhdr}

% Comments 
\usepackage{comment, chngpage, appendix, verbatim}

% Bib setup
\usepackage[backend=biber, bibencoding=utf8, sorting=none,]{biblatex}
\addbibresource{references.bib}

% Used for code snippets
\usepackage{listings}

% Use the palatino font
\usepackage[sc]{mathpazo}
\linespread{1.1} % set line spread for palatino font

% Used to create hyperrefs, \href, bliver brugt meget til links man gerne vil have har en anden tekst
\usepackage[hidelinks]{hyperref}
\hypersetup{%
	%pdfpagelabels=true,%
	plainpages=false,%
	pdfauthor={Author(s)},%
	pdftitle={Title},%
	pdfsubject={Subject},%
	bookmarksnumbered=true,%
	colorlinks=false,%
	citecolor=black,%
	filecolor=black,%
	linkcolor=black,% you should probably change this to black before printing
	urlcolor=black,%
	pdfstartview=FitH%
}

% Tables
% Make tables able to go into multible rows
\usepackage{multirow}

% Make todo tables 
\setlength{\marginparwidth}{2cm} % make sure todo tables fit
\usepackage{todonotes}

% an custom code snippet in latex
\usepackage{assets/blockcoding}

